\chapter{Function Words and Particles}\label{ch:function-words}

\section{Overview}

Function words in Iridian form several small, largely closed classes whose
importance is out of proportion to their size. Unlike nouns and verbs, they
contribute little lexical content of their own; instead they organize the
clause, mark relations between constituents, and encode the speaker's stance
toward what is being said. Among them, the preverbal particle series is
especially central, since it interacts directly with predicate structure,
aspect, interrogation, and discourse framing.

\section{Adverbial particles}\label{sec:adv-particles}

\subsection{In general}\label{sec:adv-particles-general}

Adverbial particles are a small series of words that are phonologically
proclitic and syntactically tied to the left edge of the predicate domain. In
their unmarked position they stand immediately before the predicate they scope
over: before a finite verb in verbal clauses, before the existential verb in
existential clauses, and before the predicate nominal in zero-copula
predications. Their proclitic phonology is therefore best understood as a
surface reflex of a fixed structural position, not as the sole reason for
treating them as a class.

The traditional Iridian term for these items is \ird{zemo} `grain, particle'.
The label \term{adverbial particle} is retained here because many of them behave
broadly like clause adverbs, but the class is wider than ordinary manner or
temporal adverbs. Particles may contribute phasal contour, interrogation,
addition, affirmation, reportative force, evidential stance, mirativity, and
other clause-level meanings. What unifies them is not a single semantic type
but a shared positional behavior and a fixed hierarchy of scope.

\begin{table}[ht]
	\footnotesize\sffamily
	\caption{Adverbial particle classes.}\label{tab:adv-particle-classes}
	\medskip
	\begin{tabularx}{\textwidth}{l L L l}
		\toprule
		{\sc class} & {\sc members} & {\sc core function and scope} & {\sc position} \\
		\midrule
		Class 1 &
		\ird{v}, \ird{s} &
		event-internal or predicational phasal contour; closest scope over the event or state itself &
		closest to predicate \\
		Class 2 &
		\ird{by}, \ird{och}, \ird{moc}, \ird{či}, \ird{kamo}/\ird{kam}, \ird{dlí}, \ird{jedy} &
		propositional and illocutionary meanings such as interrogation, addition, affirmation, reportative force, exclusivity, anteriority, and contrast &
		middle \\
		Class 3 &
		\ird{oče}, \ird{nadě}, \ird{hlavdy}, \ird{ktady}, \ird{mlada} &
		speaker stance over the whole utterance, including mirative, consequential, inferential, speculative, and hypothetical meanings &
		farthest from predicate \\
		\bottomrule
	\end{tabularx}
\end{table}

The basic ordering of the series is thus as follows:

\pex
\a \ird{Class 3 > Class 2 > Class 1 > predicate}
\xe

\noindent This ordering is best explained semantically: lower particles take
narrower scope and higher particles take scope over progressively larger units
of meaning. The familiar phonological tendency for shorter particles to stand
closer to the predicate and longer ones farther away is real, but it is better
treated as secondary evidence of grammaticalization depth than as the primary
source of the ordering itself.

The unmarked placement of all three classes is within this preverbal series.
The highest particles may, in marked discourse use, be placed farther left as
sentence-framing elements; this possibility is most plausible for the Class~3
particles and is taken up when those particles are discussed individually. In
the ordinary clause, however, the default pattern is strictly preverbal.

\pex
\a
\begingl
\gla Janek v piašček.//
\glb Janek.\Dir{} already eat-\Antip{}-\Pf{}//
\glft \trsl{Janek has already eaten.}//
\endgl
\a
\begingl
\gla Tak tóm Janku oče by v koupnik?//
\glb this book Janek.\Trs{} mirative question already buy-\Trs{}-\Pf{}//
\glft \trsl{Has Janek, surprisingly, already bought this book?}//
\endgl
\a
\begingl
\gla Marek v doktor.//
\glb Marek.\Dir{} already doctor.\Dir{}//
\glft \trsl{Marek is already a doctor.}//
\endgl
\xe

\subsection{Class 1: \ird{v} and \ird{s}}\label{sec:class1-particles}

The lowest slot in the particle sequence is occupied by either \ird{v} or
\ird{s}, never by both at once. These are best analysed as phasal particles
anchored to a reference time, not as direct exponents of the aspect-alignment
suffixes described in Section~\ref{sec:aspect}. Broadly speaking, \ird{v}
signals that the relevant state or event has already been attained, or is taken
to be close to attainment, by the reference point. \ird{s}, by contrast,
presents the relevant interval as still open, still continuing, or not yet
exhausted at that same point.

Because these meanings interact with viewpoint aspect, \ird{v} tends to occur
most naturally with bounded, attained, or resultative predications, while
\ird{s} prefers ongoing, persistent, or still-pending predications. The
correlation is strong, but it is not absolute. Class~1 particles do not replace
aspect morphology; rather, they add a phasal interpretation to the aspectual
viewpoint supplied elsewhere in the predicate.

\noindent\textit{Morphophonological note.} Before a following word beginning
with \ird{v}, \ird{v} surfaces as \ird{ve}; before a following word beginning
with a sibilant, \ird{s} surfaces as \ird{se}.

\subsubsection{Core meaning relative to a reference time}

With present-time reference, the contrast between \ird{v} and \ird{s} is often
most transparent with stative or nominal predicates. \ird{V} marks the state as
already reached; \ird{s} presents it as continuing.

\pex
\a
\begingl
\gla Marek doktor.//
\glb Marek.\Dir{} doctor.\Dir{}//
\glft \trsl{Marek is a doctor.}//
\endgl
\a
\begingl
\gla Marek v doktor.//
\glb Marek.\Dir{} already doctor.\Dir{}//
\glft \trsl{Marek is already a doctor.}//
\endgl
\a
\begingl
\gla Marek s doktor.//
\glb Marek.\Dir{} still doctor.\Dir{}//
\glft \trsl{Marek is still a doctor.}//
\endgl
\xe

The same placement rule holds with existential predicates. Here \ird{s}
indicates persistence, while \ird{v} tends to mark a newly attained or already
established situation.

\pex
\a
\begingl
\gla Marek dumě na jeca.//
\glb Marek.\Dir{} house.\Ind{} in \Exst{}-\Cont{}//
\glft \trsl{Marek is at home.}//
\endgl
\a
\begingl
\gla Marek dumě na s jeca.//
\glb Marek.\Dir{} house.\Ind{} in still \Exst{}-\Cont{}//
\glft \trsl{Marek is still at home.}//
\endgl
\xe

\subsubsection{With explicit temporal adjuncts}

When the clause contains an explicit time expression, the reference point
introduced by that expression becomes central to the interpretation of
\ird{v} and \ird{s}. With future-oriented predicates, \ird{v} presents the
event as already close from the speaker's current vantage point, or as one that
will have been reached by the named time; \ird{s} presents the interval leading
up to that time as still open and not yet exhausted.

\pex
\a
\begingl
\gla Janek sobotě ščenách.//
\glb Janek Saturday.\Ind{} arrive-\Ctp{}//
\glft \trsl{Janek will arrive on Saturday.}//
\endgl
\a
\begingl
\gla Janek sobotě v ščenách.//
\glb Janek Saturday.\Ind{} already arrive-\Ctp{}//
\glft \trsl{Janek will already have arrived by Saturday / Saturday already feels close for Janek's arrival.}//
\endgl
\a
\begingl
\gla Janek sobotě se ščenách.//
\glb Janek Saturday.\Ind{} still arrive-\Ctp{}//
\glft \trsl{Janek will arrive on Saturday, and there is still time before then.}//
\endgl
\xe

With past-oriented predicates, \ird{s} may likewise measure the interval
between the event time and the present, yielding the sense that a considerable
stretch of time has elapsed since that event. In the same environment,
\ird{v} usually retains its ordinary already-attained reading instead of
contributing the same interval effect.

\pex
\a
\begingl
\gla Janek sobotě ščenik.//
\glb Janek Saturday.\Ind{} arrive-\Pf{}//
\glft \trsl{Janek arrived on Saturday.}//
\endgl
\a
\begingl
\gla Janek sobotě se ščenik.//
\glb Janek Saturday.\Ind{} still arrive-\Pf{}//
\glft \trsl{Janek arrived on Saturday, and it has been quite some time since then.}//
\endgl
\xe

\subsubsection{Interaction with aspect morphology}

The relation between Class~1 particles and aspect is systematic, but it is not
a matter of simple one-to-one matching. \ird{V} is especially natural with
predicates that present a bounded result, a completed event, or an attained
state: perfectives, retrospectives, zero-copula predications, and some
contemplatives. \ird{S} is especially natural with predicates that present an
ongoing state, an open interval, or a still-relevant course of action:
continuous and progressive predicates, persistent existential clauses, and some
contemplatives.

This means that \ird{v} overlaps partly with the retrospective, but the two are
not equivalent. The retrospective contributes prior-event perspective and
continuing relevance; \ird{v} adds the scalar meaning that the relevant result
has already been reached by the reference point. A perfective clause with
\ird{v} therefore often resembles an English \trsl{already}-perfect, while a
retrospective clause with \ird{v} places that already-attained result within a
broader frame of prior relevance.

\pex
\a
\begingl
\gla Janek ščenik.//
\glb Janek.\Dir{} arrive-\Pf{}//
\glft \trsl{Janek arrived.}//
\endgl
\a
\begingl
\gla Janek v ščenik.//
\glb Janek.\Dir{} already arrive-\Pf{}//
\glft \trsl{Janek has already arrived.}//
\endgl
\a
\begingl
\gla Pražě de stožaní.//
\glb Prague.\Ind{} to go-\Ret{}//
\glft \trsl{(I) have been to Prague before.}//
\endgl
\xe

\noindent In a framed retrospective clause, \ird{v} shows that the relevant
prior event had already been completed by the framing event:

\pex
\begingl
\gla Marek ščenu Janek v piaščaní.//
\glb Marek.\Dir{} arrive-\Seq{} Janek.\Dir{} already eat-\Antip{}-\Ret{}//
\glft \trsl{When Marek arrived, Janek had already eaten.}//
\endgl
\xe

With contemplatives, the same contrast appears in a different form: \ird{v}
often yields an imminent reading (\trsl{already about to}), whereas \ird{s}
preserves the sense that the projected event or intention is still in force
(\trsl{still going to}, \trsl{still intending to}).

\subsubsection{Negated predicates}

Under negation, the core meanings of \ird{v} and \ird{s} remain visible as a
scope interaction. \ird{V} with a negative predicate usually marks cessation or
loss of a previously attained state, yielding translations such as
\trsl{no longer} or \trsl{no more}. \ird{S} with a negative predicate marks
non-attainment, delay, or an interval that remains open, yielding
\trsl{not yet}. Stronger emphatic readings are possible in context, but these
follow from the basic contrast rather than replacing it.

\pex
\a
\begingl
\gla Pivo v zada.//
\glb beer.\Dir{} already \NegExst{}-\Cont{}//
\glft \trsl{There is no more beer.}//
\endgl
\a
\begingl
\gla Pivo s zada.//
\glb beer.\Dir{} still \NegExst{}-\Cont{}//
\glft \trsl{There is no beer yet.}//
\endgl
\a
\begingl
\gla Marek doktor v staly.//
\glb Marek.\Dir{} doctor.\Dir{} already \NegCop{}//
\glft \trsl{Marek is no longer a doctor.}//
\endgl
\a
\begingl
\gla Marek doktor s staly.//
\glb Marek.\Dir{} doctor.\Dir{} still \NegCop{}//
\glft \trsl{Marek is not yet a doctor.}//
\endgl
\xe

\subsubsection{Extended uses of \ird{s}}

Besides its core continuative meaning, \ird{s} has several extended uses that
are best understood as developments from the idea of an interval that remains
open or of an additional step on a scale. With perfective predicates,
especially in narratives of inconvenience, effort, or unexpected inclusion,
\ird{s} may be translated as \trsl{even}. This use often carries annoyance,
commiseration, or rhetorical emphasis; the affirmative particle \ird{moc} may
reinforce the interpretation, though the details belong properly to the Class~2
discussion.

\pex
\a
\begingl
\gla Magazinam de bych s stojnik!//
\glb store.\Ind{} into yesterday still go-\Pf{}//
\glft \trsl{I even went to the store yesterday!}//
\endgl
\a
\begingl
\gla Marek Štěpánovímně sprivam moc se stojnik.//
\glb Marek.\Dir{} Štěpánov neighborhood.\Ind{} toward affirmative still go-\Pf{}//
\glft \trsl{Marek even went to Štěpánov.}//
\endgl
\xe

In questions, the same extension yields \trsl{else} or \trsl{in addition}:

\pex
\begingl
\gla Mám Marku s piaštnik?//
\glb what Marek.\Trs{} still eat-\Trs{}-\Pf{}//
\glft \trsl{What else did Marek eat?}//
\endgl
\xe

These meanings should not be treated as the core semantics of \ird{s}. They are
contextually licensed extensions of the general continuative/additive profile of
the particle.

\subsubsection{Yes-no questions}

The yes-no question particle is \ird{by}; \ird{v} and \ird{s} do not themselves
form questions. Rather, they specify whether the speaker asks about an already
attained result, an ongoing situation, or an interval that remains open. This
distinction becomes especially clear when the question concerns completion or
timing relative to a reference point.

\begin{table}[ht]
\footnotesize\sffamily
\caption{Class 1 particles in yes-no questions.}\label{tab:v-s-questions}
\medskip
\begin{tabularx}{\textwidth}{l L L}
\toprule
{\sc aspect} & {\sc with \ird{v}} & {\sc with \ird{s}} \\
\midrule
Perfective &
\ird{Janek by v piašček?} \newline
\trsl{Has Janek already eaten?} &
\ird{Janek by s piašček?} \newline
\trsl{Did Janek still manage to eat?} \\
Progressive &
\ird{Janek by v piaštime?} \newline
\trsl{Is Janek already eating?} &
\ird{Janek by s piaštime?} \newline
\trsl{Is Janek still eating?} \\
Contemplative &
\ird{Janek by v piaštách?} \newline
\trsl{Is Janek already about to eat?} &
\ird{Janek by s piaštách?} \newline
\trsl{Is Janek still planning to eat?} \\
\bottomrule
\end{tabularx}
\end{table}

\noindent The general syntax of yes-no questions is discussed in
Section~\ref{sec:questions} of Chapter~\ref{ch:simple}; what matters here is
the contribution of the Class~1 particle to the temporal-phasal profile of the
question.

\section{Connectives and conjunctions}

The remaining function-word classes are treated in later sections of this
chapter.

\section{Discourse-connective particles}

Many discourse-connective particles arise historically from older predicative or
clitic material. Their interaction with the sentence-level particle system will
be described below.

\section{Interjections}

Interjections stand outside ordinary clause structure and are not part of the
preverbal particle sequence described in Section~\ref{sec:adv-particles}.
